\documentclass[runningheads]{llncs}

% ---------------------------------------------------------------
% Include basic ECCV package
 
% TODO REVIEW: Insert your submission number below by replacing '*****'
% TODO FINAL: Comment out the following line for the camera-ready version
\usepackage[review,year=2026,ID=*****]{eccv}
% TODO FINAL: Un-comment the following line for the camera-ready version
%\usepackage{eccv}

% OPTIONAL: Un-comment the following line for a version which is easier to read
% on small portrait-orientation screens (e.g., mobile phones, or beside other windows)
%\usepackage[mobile]{eccv}


% ---------------------------------------------------------------
% Other packages

% Commonly used abbreviations (\eg, \ie, \etc, \cf, \etal, etc.)
\usepackage{eccvabbrv}

% Include other packages here, before hyperref.
\usepackage{graphicx}
\usepackage{booktabs}

% The "axessiblity" package can be found at: https://ctan.org/pkg/axessibility?lang=en
\usepackage[accsupp]{axessibility}  % Improves PDF readability for those with disabilities.


% ---------------------------------------------------------------
% Hyperref package

% It is strongly recommended to use hyperref, especially for the review version.
% Please disable hyperref *only* if you encounter grave issues.
% hyperref with option pagebackref eases the reviewers' job, but should be disabled for the final version.
%
% If you comment hyperref and then uncomment it, you should delete
% main.aux before re-running LaTeX.
% (Or just hit 'q' on the first LaTeX run, let it finish, and you
%  should be clear).

% TODO FINAL: Comment out the following line for the camera-ready version
%\usepackage[pagebackref,breaklinks,colorlinks,citecolor=eccvblue]{hyperref}
% TODO FINAL: Un-comment the following line for the camera-ready version
\usepackage{hyperref}

% Support for ORCID icon
\usepackage{orcidlink}


\begin{document}

% ---------------------------------------------------------------
% TODO REVIEW: Replace with your title
\title{Recovering Metric Scale in Pretrained Image-to-3D Generators Without Backbone Retraining}

% TODO REVIEW: If the paper title is too long for the running head, you can set
% an abbreviated paper title here. If not, comment out.
\titlerunning{Metric Scale Recovery for Image-to-3D Generators}

% TODO FINAL: Replace with your author list. 
% Include the authors' OCRID for the camera-ready version, if at all possible.
\author{First Author\inst{1}\orcidlink{0000-1111-2222-3333} \and
Second Author\inst{2,3}\orcidlink{1111-2222-3333-4444} \and
Third Author\inst{3}\orcidlink{2222--3333-4444-5555}}

% TODO FINAL: Replace with an abbreviated list of authors.
\authorrunning{F.~Author et al.}
% First names are abbreviated in the running head.
% If there are more than two authors, 'et al.' is used.

% TODO FINAL: Replace with your institution list.
\institute{Princeton University, Princeton NJ 08544, USA \and
Springer Heidelberg, Tiergartenstr.~17, 69121 Heidelberg, Germany
\email{lncs@springer.com}\\
\url{http://www.springer.com/gp/computer-science/lncs} \and
ABC Institute, Rupert-Karls-University Heidelberg, Heidelberg, Germany\\
\email{\{abc,lncs\}@uni-heidelberg.de}}

\maketitle


\begin{abstract}
  The abstract should concisely summarize the contents of the paper. 
  While there is no fixed length restriction for the abstract, it is recommended to limit your abstract to approximately 150 words.
  Please include keywords as in the example below. 
  This is required for papers in LNCS proceedings.
  \keywords{First keyword \and Second keyword \and Third keyword}
\end{abstract}


\section{Introduction}
\label{sec:intro}

Image-to-3D generative models have reached a level of visual fidelity where a single photograph
is sufficient to produce a detailed mesh or Gaussian splat of an object with plausible geometry
and texture. Foundations such as SAM 3D Objects~\cite{sam3d} and TRELLIS~\cite{trellis} train
on millions of synthetic and real objects, learning rich priors over 3D shape and surface
appearance. Yet these models share a fundamental blind spot: their outputs live in a canonical
unit cube, $[-0.5, 0.5]^3$, with no connection to physical units. A reconstructed bowl and a
reconstructed car are indistinguishable in scale.

This missing dimension --- metric scale --- is not a cosmetic limitation. Robotic manipulation
requires knowing whether a grasped object is five centimetres or fifty. Augmented reality demands
that a reconstructed mug placed on a physical table displaces a physically plausible volume.
Scene reconstruction from multiple objects depends on consistent metric embedding. Without
absolute scale, a 3D generative model is an artist's tool, not a spatial computing primitive.

\paragraph{Why metric scale is hard.}
Recovering metric scale from a single image is fundamentally under-constrained without a prior
on object size. The standard workarounds each carry a cost. \emph{Category-mean priors} assume
the object is an average specimen of its class --- a heuristic that incurs $\sim$10.9\% mean
absolute percentage error (MAPE) on NOCS-Real275~\cite{nocs} even when the category is known.
\emph{Separate monocular metric depth models}~\cite{metric3d,zoedepth} can estimate depth in
physical units but require an additional inference pass, introduce a second set of parameters,
and still require a geometric bridge between pixel-depth and object volume.
\emph{Backbone retraining} --- modifying the 3D generator itself to reason about physical scale
--- demands compute budgets exceeding 100K GPU-hours for models of this class, which would also
destroy the generalisation properties built into the pretrained weights.

\paragraph{The hidden anchor.}
We observe that modern image-to-3D pipelines based on monocular depth estimation already contain
the metric signal needed for scale recovery --- it is simply being discarded. SAM 3D Objects uses
MoGe~\cite{moge} to lift the input image into a metric pointmap: a dense field of $(x, y, z)$
coordinates in metres aligned to the camera frame. From this pointmap the pipeline extracts a
global scale statistic (\texttt{pointmap\_scale}) and a vertical offset
(\texttt{pointmap\_shift\_z}) to normalise input conditioning --- and then never uses the
original metric values again. The object's true physical extent, encoded in those pointmap
statistics and in the spatial distribution of pointmap tokens, is present at inference time
but plays no role in the output.

\paragraph{Our approach.}
We propose a lightweight recipe for injecting metric awareness into a frozen image-to-3D
generator without retraining any backbone component. A compact \emph{MetricScaleHead} fuses
MoGe pointmap statistics with pooled sparse-structure (SS) latent features into a
1024-dimensional scale token. This token is injected into the spatial latent (SLAT)
flow-matching decoder via cross-attention --- the only part of the generative model that is
fine-tuned ($\sim$100M of more than one billion total parameters), with an fp32 upcast recipe
that prevents the bf16 attention overflow we encountered in early experiments. Separately, a
differentiable \emph{aspect-ratio loss} on the SS decoder's soft voxel occupancy supervises
object shape proportions directly, closing the gap between the metric head's per-axis scale
predictions and the mesh's actual bounding-box dimensions.

\paragraph{Results.}
Trained on the real-world NOCS-Real275 subset of OmniNOCS~\cite{omninocs} and evaluated on a
held-out scene split, our method achieves \textbf{1.23\% MAPE} (0.20\,cm mean absolute error)
--- an order of magnitude below the category-prior baseline and competitive with specialised
monocular metric depth methods that require a separate inference pass. Extended to
mixed-source training on Objectron~\cite{objectron} and ARKitScenes~\cite{arkitscenes}, the
model predicts metric scale across object categories, viewpoints, and indoor environments
without sacrificing per-category accuracy on the original evaluation set.

\paragraph{Contributions.} We contribute:
\begin{itemize}
  \item The observation that metric pointmap statistics are a latent anchor already present in
        image-to-3D pipelines and amenable to lightweight supervision.
  \item A practical recipe --- MetricScaleHead, scale-token cross-attention injection, and soft
        aspect-ratio supervision --- that adds metric scale to a frozen 3D generator with
        $\sim$170M trainable parameters and no backbone modification.
  \item A training stability protocol (fp32 cross-attention upcast, linear lr warmup) that
        makes fine-tuning large sparse flow-matching decoders viable without NaN instabilities.
  \item Empirical validation across three real-world object datasets demonstrating sub-2\%
        MAPE with generalisation to unseen object categories.
\end{itemize}

\section{Related Work}
\label{sec:related}

\paragraph{Image-to-3D generative models.}
Large-scale image-to-3D systems have converged on a two-stage latent flow-matching design. TRELLIS~\cite{trellis}
introduces the Structured 3D Latent (SLAT) representation, a sparse voxel grid carrying DINOv2 visual features
extracted from multi-view renders; a rectified-flow transformer decodes SLAT tokens conditioned on a single image,
with separate VAE heads producing Gaussian splats or textured meshes via FlexiCubes. TRELLIS v2~\cite{trellis2}
replaces the multiview-supervised features with native O-Voxel representations --- a ``field-free'' sparse structure
that handles open, non-manifold, and fully-enclosed surfaces --- achieving a $16\times$ spatial compression rate
and adding explicit PBR material support across four billion parameters. SAM 3D~\cite{sam3d} extends the SLAT
backbone to natural images with significant clutter and occlusion: a 1.2B-parameter Mixture-of-Transformers
geometry model jointly estimates coarse shape and object layout, followed by a 600M-parameter sparse latent
refinement model, all trained through a multi-stage synthetic-to-real flywheel aligned to human preference.
Despite their differences, all three systems output objects normalised to a unit cube with no connection to
physical dimensions. Our work keeps these frozen backbones intact and recovers the discarded metric scale
via a lightweight head trained on real-world object datasets.

\paragraph{Monocular geometry and metric depth estimation.}
Recovering scene geometry in physical units from a single image has seen rapid recent progress.
ZoeDepth~\cite{zoedepth} fine-tunes a relative-depth backbone with per-domain metric bins, enabling
zero-shot metric depth across indoor and outdoor scenes.
Metric3D~\cite{metric3d} achieves scale-aware depth through a canonical camera normalization module that
reconciles training on datasets with heterogeneous camera intrinsics.
MoGe~\cite{moge} reframes the problem as \emph{affine-invariant} pointmap prediction: it estimates a dense
$(x,y,z)$ field using a robust optimal-alignment loss, achieving state-of-the-art relative geometry accuracy
while deliberately leaving global scale and shift unconstrained. MoGe-2~\cite{moge2} extends this by
appending a lightweight MLP head on the ViT CLS token to predict a global metric scale factor, enabling
accurate metric geometry without sacrificing relative geometry quality. MoGe is integrated as a front-end in
both SAM 3D and the SAM 3D Objects pipeline we build on: the pipeline computes alignment statistics
(\texttt{pointmap\_scale}, \texttt{pointmap\_shift\_z}) to map MoGe's output into a unit-cube conditioning
space --- and then discards these normalization factors before the generative decoder runs. Our key
observation is that these statistics are correlated with the object's true physical extent, and routing them
directly into a supervised scale head recovers metric dimensions without any additional inference pass.

\paragraph{Category-level size and pose estimation.}
NOCS~\cite{nocs} defines the normalised object coordinate space framework for recovering 6-DoF pose and
metric bounding-box dimensions for known object categories; the accompanying NOCS-Real275 benchmark
(6 categories, 7k images) is our primary evaluation set.
OmniNOCS~\cite{omninocs} substantially extends the NOCS paradigm, unifying canonical 3D bounding-box
annotations across ten data sources --- including NOCS-Real275, Objectron~\cite{objectron},
ARKitScenes~\cite{arkitscenes}, KITTI, nuScenes, and Hypersim --- to produce a dataset of 97 object
categories and 380k images, twenty times the class coverage and 200 times the instance count of
prior NOCS datasets. The accompanying NOCSformer model predicts NOCS coordinates, instance masks,
and metric 3D boxes from 2D detections via a class-agnostic ViT backbone, generalising to unseen datasets
without class-specific heads. All category-level methods share a common limitation: metric scale recovery
is implicitly conditioned on category-level size priors encoded during training, which incurs a baseline MAPE
of $\sim$10.9\% on NOCS-Real275 even with ground-truth category labels. Our method uses no category
labels at inference time --- scale is inferred from pointmap normalization statistics and sparse latent
features that are category-agnostic --- and is trained on the same OmniNOCS sources.

\paragraph{Scale recovery in frozen generative models.}
Adapter-style fine-tuning has proven effective for injecting new capabilities into large frozen models in both
language~\cite{} and vision domains~\cite{}, but the objective has typically been semantic (e.g., style, content
control) rather than metric. In 3D, methods such as Zero123, InstantMesh, and CRM recover geometry in
normalised space; recovering metric scale from their outputs requires an independent calibration step.
To our knowledge, no prior work injects a metric supervision signal into the latent decoder of a large image-to-3D
generative model while preserving its frozen backbone weights. Our MetricScaleHead and cross-attention injection
recipe fill this gap.


\section{Method}
\label{sec:method}

TODO


\section{Experiments}
\label{sec:experiments}

TODO


\section{Conclusion}
\label{sec:conclusion}

TODO


\section*{Acknowledgements}
TODO

% ---- Bibliography ----
%
% BibTeX users should specify bibliography style 'splncs04'.
% References will then be sorted and formatted in the correct style.
%
\bibliographystyle{splncs04}
\bibliography{main}
\end{document}
